\section{Лекция 1. Основы дискретной математики}
\noeqref{eq:subset_def, eq:set_eq_def, eq:strict_subset_def, eq:concat_id}
\noeqref{eq:concat_assoc, eq:cart_prod_n, eq:cart_pow_n, eq:cart_pow_omega}
\noeqref{eq:kleene_star, eq:kleene_plus, eq:cantor_main}

\subsection*{Примерная структура курса}

\begin{enumerate}
    \item Множества и отношения \hfill $\sim 2$ лекции
    \item Комбинаторика \hfill $\sim 4$ лекции
    \item Теория графов \hfill $\sim 4$ лекции
    \item Теория кодирования \hfill $\sim 2$ лекции
    \item Булевы функции \hfill $\sim$ остальное время
\end{enumerate}

\subsection{Множества (Наивная теория множеств)}

\begin{definition}
    \textdef{Множество} --- совокупность элементов, мыслимая как единое 
    целое.
\end{definition}

\begin{designation} ~
    \begin{align*}
        X \quad &\text{--- множества обозначаются заглавными буквами;} \\
        x \quad &\text{--- элементы множества обозначаются строчными буквами;} \\
        x \in X \quad &\text{--- элемент } x \text{ принадлежит множеству } X; \\
        x \notin X \quad &\text{--- элемент } x \text{ не принадлежит множеству } X.
    \end{align*}
\end{designation}

\begin{remark}[Способы задания и специальные множества] ~
    Способы задания множеств:
    \begin{itemize}
        \item \textbf{Перечислением элементов:} $\{x_1, x_2, \dots, x_n\}$ 
        (применимо для конечных множеств).
        \item \textbf{Характеристическим свойством:} $\{x \mid \phi(x)\}$ 
        или $\{x \in X \mid \phi(x)\}$, где $\phi(x)$ --- некоторое условие.
    \end{itemize}

    Специальные множества, используемые в курсе:
    \begin{itemize}
        \item \textdef{Булево множество:} $\B = \{0, 1\}$;
        \item \textdef{Числовые множества:} $\N, \Z, \Q, \R, \C$.
    \end{itemize}
\end{remark}

\subsubsection{Операции над множествами}

\begin{definition}
    Пусть $X$ и $Y$ --- некоторые множества. Определим основные 
    теоретико-множественные операции над ними:
    \begin{enumerate}
        \item \textdef{Объединение:} 
        $X \cup Y = \{x \mid x \in X \text{ или } x \in Y\}$;
        \item \textdef{Пересечение:} 
        $X \cap Y = \{x \mid x \in X \text{ и } x \in Y\}$;
        \item \textdef{Разность:} 
        $X \setminus Y = \{x \mid x \in X \text{ и } x \notin Y\}$;
        \item \textdef{Симметрическая разность:} 
        $X \triangle Y = (X \setminus Y) \cup (Y \setminus X) = 
        (X \cup Y) \setminus (X \cap Y)$;
        \item \textdef{Дополнение:} 
        $C_X Y = X \setminus Y$ (дополнение множества $Y$ до универсального 
        множества $X$, где $Y \subseteq X$).
    \end{enumerate}
\end{definition}

\subsubsection{Подмножества и равенство множеств}

\begin{definition}
    Множество $X$ называется \textdef{подмножеством} множества $Y$ 
    (обозначается $X \subseteq Y$), если для любого элемента $x$ верна 
    импликация:
    \begin{equation} \label{eq:subset_def}
        (x \in X) \implies (x \in Y) \tag{1}
    \end{equation}
\end{definition}

\begin{definition}
    Множества $X$ и $Y$ называются \textdef{равными} ($X = Y$), если они 
    состоят из одних и тех же элементов, то есть выполняется взаимное 
    включение:
    \begin{equation} \label{eq:set_eq_def}
        X = Y \iff (X \subseteq Y \quad \text{и} \quad Y \subseteq X) \tag{2}
    \end{equation}
\end{definition}

\begin{definition}
    Множество $X$ называется \textdef{строгим (собственным) подмножеством} 
    множества $Y$ (обозначается $X \subset Y$), если:
    \begin{equation} \label{eq:strict_subset_def}
        X \subset Y \iff (X \subseteq Y \quad \text{и} \quad X \neq Y) \tag{3}
    \end{equation}
\end{definition}

\begin{definition}
    \textdef{Пустое множество} $\varnothing$ --- это уникальное множество, 
    не содержащее ни одного элемента.
\end{definition}

\begin{remark}
    Пустое множество является подмножеством любого множества: 
    $\forall X : \varnothing \subseteq X$.
\end{remark}

\begin{remark}
    Следует строго различать пустое множество $\varnothing$ и множество 
    $\{\varnothing\}$, состоящее из пустого множества. Множество 
    $\varnothing$ не содержит элементов ($\abs{\varnothing} = 0$), тогда как 
    $\{\varnothing\}$ содержит ровно один элемент --- само $\varnothing$ 
    ($\abs{\{\varnothing\}} = 1$). Поэтому $\varnothing \neq \{\varnothing\}$, 
    хотя одновременно выполняются соотношения $\varnothing \in \{\varnothing\}$ 
    и $\varnothing \subseteq \{\varnothing\}$.
\end{remark}

\subsubsection{Кортежи, слова и декартово произведение}

\begin{definition}
    \textdef{Совокупность} --- первичное, неопределяемое строго понятие, 
    обозначающее мысленное объединение объектов в единое целое.
\end{definition}

\begin{definition}
    \textdef{Упорядоченная пара} $(x, y)$ --- двухэлементная совокупность, 
    в которой зафиксирован порядок следования элементов: $(x, y) = (u, v)$ 
    тогда и только тогда, когда $x = u$ и $y = v$.
\end{definition}

\begin{definition}
    \textdef{Кортеж} длины $n$ --- конечный упорядоченный набор из $n$ 
    элементов, обозначаемый $(x_1, x_2, \dots, x_n)$. В отличие от 
    множества, в кортеже важен порядок элементов, а сами элементы могут 
    повторяться.
\end{definition}

\begin{definition}
    \textdef{Последовательность} --- функция, областью определения которой 
    является множество натуральных чисел $\N$ (бесконечная 
    последовательность) или его начальный отрезок $\{1, 2, \dots, n\}$ 
    (конечная последовательность).
\end{definition}

\begin{definition}
    \textdef{Алфавит} --- любое конечное непустое множество, элементы 
    которого называются \textdef{символами} или \textdef{буквами}.
\end{definition}

\begin{definition}
    \textdef{Слово конечной длины} (или \textdef{строка}) --- конечная 
    последовательность (кортеж) символов, выбранных из заданного алфавита.
\end{definition}

\begin{definition}
    \textdef{Пустое слово} --- слово, не содержащее ни одного символа 
    (его длина равна $0$). Обозначается символом $\eps$ или $\Lambda$.
\end{definition}

\begin{definition}
    \textdef{Конкатенация} (сцепление) слов --- бинарная операция, 
    заключающаяся в последовательном приписывании одного слова в конец 
    другого. Если даны слова $u = a_1 a_2 \dots a_k$ и 
    $v = b_1 b_2 \dots b_m$, то их конкатенацией называется слово 
    $u v = a_1 a_2 \dots a_k b_1 b_2 \dots b_m$ длины $k + m$.
\end{definition}

\begin{properties}[Свойства конкатенации и пустого слова] ~
    \begin{enumerate}
        \item Пустое слово $\eps$ является нейтральным элементом 
        относительно операции конкатенации. Для любого слова $w$ верно:
        \begin{equation} \label{eq:concat_id}
            w \eps = \eps w = w \tag{4}
        \end{equation}
        \item Операция конкатенации ассоциативна:
        \begin{equation} \label{eq:concat_assoc}
            (u v) w = u (v w) \tag{5}
        \end{equation}
        \item В общем случае операция конкатенации некоммутативна: 
        $u v \neq v u$.
    \end{enumerate}
\end{properties}

\begin{definition}
    \textdef{Декартовым (прямым) произведением} двух множеств $X$ и $Y$ 
    называется множество всех упорядоченных пар, в которых первый элемент 
    принадлежит $X$, а второй --- $Y$:
    \[
        X \times Y = \{(x, y) \mid x \in X, \ y \in Y\}
    \]
    Обобщением на $n$ множеств ($n \in \N$) является множество кортежей 
    длины $n$:
    \begin{equation} \label{eq:cart_prod_n}
        X_1 \times X_2 \times \dots \times X_n = 
        \{\underbrace{(x_1, \dots, x_n)}_{\text{кортеж}} \mid 
        x_i \in X_i, \ i = \overline{1, n}\} \tag{6}
    \end{equation}
\end{definition}

\begin{remark}
    Строго говоря, операция декартова произведения не ассоциативна:
    \begin{align*}
        X_1 \times X_2 \times X_3 &= \{(x_1, x_2, x_3)\} 
        \quad \text{(множество упорядоченных троек)}, \\
        (X_1 \times X_2) \times X_3 &= \{((x_1, x_2), x_3)\} 
        \quad \text{(множество пар, где первый элемент --- пара)}.
    \end{align*}
    Однако между этими множествами существует естественная биекция, поэтому 
    на практике их часто отождествляют.
\end{remark}

\begin{definition}
    \textdef{$n$-й декартовой степенью} множества $X$ ($n \in \N$) 
    называется декартово произведение $n$ копий множества $X$:
    \begin{equation} \label{eq:cart_pow_n}
        X^n = \underbrace{X \times X \times \dots \times X}_{n \text{ раз}} = 
        \{\underbrace{(x_1, \dots, x_n)}_{\text{кортеж}} \mid 
        x_i \in X, \ i = \overline{1, n}\} \tag{7}
    \end{equation}
    По соглашению \textdef{нулевая декартова степень} $X^0$ содержит 
    единственный элемент --- пустой кортеж (пустое слово):
    \[
        X^0 = \{ () \} = \{ \eps \}
    \]
\end{definition}

\begin{definition}
    Пусть символ $\omega$ обозначает счётную бесконечность. 
    \textdef{Счётной декартовой степенью} множества $X$ называется 
    множество всех бесконечных последовательностей ($\omega$-слов) 
    из элементов $X$:
    \begin{equation} \label{eq:cart_pow_omega}
        X^\omega = \left\{
            \underbrace{(x_1, x_2, \dots, x_n, \dots)}_{
                \text{последовательность или } \omega\text{-слово}
            } \;\middle|\; x_i \in X, \ i = \overline{1, \infty}
        \right\} \tag{8}
    \end{equation}
\end{definition}

\begin{designation} ~
    \begin{itemize}
        \item Множество всех слов конечной длины над множеством $X$ 
        (\textdef{звезда Клини} или итерация Клини):
        \begin{equation} \label{eq:kleene_star}
            X^* = \bigcup_{k=0}^{\infty} X^k \tag{9}
        \end{equation}
        \item Множество всех непустых слов конечной длины (\textdef{плюс 
        Клини}):
        \begin{equation} \label{eq:kleene_plus}
            X^+ = \bigcup_{k=1}^{\infty} X^k \tag{10}
        \end{equation}
        \item Множество всех конечных и бесконечных слов:
        \[
            X^{\le \omega} = X^* \cup X^\omega
        \]
    \end{itemize}
\end{designation}

\subsection{Мощность множеств}

\subsubsection{Счётные множества и кардинальные числа}

\begin{definition}
    Множества $X$ и $Y$ называются \textdef{равномощными} (обозначается 
    $X \sim Y$), если между ними существует взаимно однозначное 
    соответствие (биекция).
\end{definition}

\begin{definition}
    \textdef{Мощностью конечного множества} $X$ называется количество 
    элементов в нём. Обозначается $\abs{X}$.
\end{definition}

\begin{definition}
    \textdef{Кардинальное число} (или \textdef{мощность}) --- это 
    характеристика множества, обобщающая понятие количества элементов на 
    бесконечные множества: два множества имеют одинаковое кардинальное 
    число тогда и только тогда, когда они равномощны. Для конечных множеств 
    кардинальные числа совпадают с натуральными числами и нулём: 
    $0, 1, 2, \dots$. Для бесконечных множеств используются специальные 
    символы (\textdef{алефы}).
\end{definition}

\begin{definition}
    Множество $X$ называется \textdef{счётным}, если оно равномощно 
    множеству натуральных чисел $\N$ ($X \sim \N$).
\end{definition}

\begin{definition}
    Множество называется \textdef{не более чем счётным}, если оно либо 
    конечно, либо счётно.
\end{definition}

\begin{designation} ~
    Стандартные обозначения бесконечных мощностей:
    \begin{align*}
        \abs{\N} = \abs{\Z} = \abs{\Q} &= \aleph_0 \quad \text{(алеф-нуль)}, \\
        \abs{\R} &= \mathfrak{c} \quad \text{(мощность континуума)}.
    \end{align*}
    Упорядоченная шкала кардинальных чисел имеет вид: 
    $0, 1, 2, \dots, \aleph_0, \aleph_1, \dots$
\end{designation}

\begin{theorem}[\textbf{Свойства счётных множеств}] ~
    \begin{enumerate}
        \item Объединение конечного числа конечных множеств конечно.
        \item Объединение счётного числа конечных множеств не более чем 
        счётно.
        \item Объединение конечного числа счётных множеств счётно.
        \item Объединение счётного числа счётных множеств счётно.
    \end{enumerate}
\end{theorem}

\begin{proof}[Доказательство свойств счётных множеств]~\\
    \begin{quote}
        Докажем последовательно каждое из четырёх утверждений теоремы.

        \textbf{1. Объединение конечного числа конечных множеств конечно.}\\
        Пусть даны конечные множества $A_1, A_2, \dots, A_n$. Так как 
        каждое из них конечно, оно содержит конечное число элементов: 
        $\abs{A_i} = k_i \in \N \cup \{0\}$. Мощность их объединения не 
        превосходит суммы их мощностей:
        \[
            \abs{\bigcup_{i=1}^n A_i} \le \sum_{i=1}^n k_i
        \]
        Сумма конечного числа целых неотрицательных чисел конечна, 
        следовательно, объединение является конечным множеством.

        \textbf{2. Объединение счётного числа конечных множеств не более 
        чем счётно.}\\
        Пусть дано счётное семейство конечных множеств $A_1, A_2, A_3, \dots$. 
        Элементы каждого конечного множества $A_i$ можно занумеровать: 
        $A_i = \{a_{i1}, a_{i2}, \dots, a_{i k_i}\}$. Выпишем элементы всех 
        множеств в единую последовательность: сначала все элементы $A_1$, 
        затем элементы $A_2$, и так далее, пропуская при этом дубликаты 
        (элементы, которые уже встречались на предыдущих шагах). Поскольку 
        каждое множество $A_i$ конечно, любой элемент из объединения рано 
        или поздно получит свой уникальный натуральный номер. Следовательно, 
        объединение конечно или счётно, то есть не более чем счётно.

        \textbf{3. Объединение конечного числа счётных множеств счётно.}\\
        Пусть дано конечное число счётных множеств $A_1, A_2, \dots, A_n$. 
        Выстроим элементы каждого множества $A_i$ в бесконечную 
        последовательность: $A_i = \{a_{i1}, a_{i2}, a_{i3}, \dots\}$. 
        Сформируем единую последовательность, выбирая по одному элементу из 
        каждого множества по очереди («по столбцам»):
        \[
            a_{11}, a_{21}, \dots, a_{n1}, \quad
            a_{12}, a_{22}, \dots, a_{n2}, \quad \dots
        \]
        Пропуская возможные дубликаты, мы сопоставим каждому уникальному 
        элементу объединения натуральный номер (без пропусков дубликатов 
        элемент $a_{ij}$ имел бы номер $(j - 1)n + i$). Так как объединение 
        содержит счётное подмножество $A_1$, оно бесконечно. Бесконечное 
        занумерованное множество является счётным.

        \textbf{4. Объединение счётного числа счётных множеств счётно.}\\
        Пусть дано счётное семейство счётных множеств $A_1, A_2, A_3, \dots$, 
        где $A_i = \{a_{i1}, a_{i2}, a_{i3}, \dots\}$. Запишем все элементы 
        в бесконечную таблицу, где $i$-я строка содержит элементы множества 
        $A_i$. Обойдём эту таблицу диагональной «змейкой» (метод Кантора), 
        последовательно перебирая диагонали с постоянной суммой индексов 
        $i + j = 2, 3, 4, \dots$:

        \begin{figure}[H]
            \centering
            \begin{tikzpicture}[>={Stealth[scale=0.8]}, scale=1.0]
                % Бесконечная матрица элементов
                \matrix (m) [
                    matrix of math nodes,
                    row sep=1.5em,
                    column sep=1.5em,
                    nodes={anchor=center}
                ] {
                    A_1 & a_{11} & a_{12} & a_{13} & \dots \\
                    A_2 & a_{21} & a_{22} & a_{23} & \dots \\
                    A_3 & a_{31} & a_{32} & a_{33} & \dots \\
                    \vdots & \vdots & \vdots & \vdots & \ddots \\
                };

                % Диагональный обход "змейкой" Кантора
                \draw[->, blue, thick] (m-1-2) -- (m-1-3);
                \draw[->, blue, thick] (m-1-3) -- (m-2-2);
                \draw[->, blue, thick] (m-2-2) -- (m-3-2);
                \draw[->, blue, thick] (m-3-2) -- (m-2-3);
                \draw[->, blue, thick] (m-2-3) -- (m-1-4);

                % Переход на следующую диагональ
                \draw[->, blue, thick] (m-1-4) -- (m-1-5);
                \draw[->, blue, thick] (m-1-5) -- (m-2-4);
                \draw[->, blue, thick] (m-2-4) -- (m-3-3);
                \draw[->, blue, thick] (m-3-3) -- (m-3-4);
            \end{tikzpicture}
            \caption{Диагональный обход счётного семейства счётных множеств 
            по методу Кантора.}
        \end{figure}

        Порядок обхода элементов имеет вид:
        \[
            a_{11}, a_{12}, a_{21}, a_{31}, a_{22}, a_{13}, a_{14}, 
            a_{23}, a_{32}, \dots
        \]
        Пропуская встречающиеся при обходе дубликаты, мы присвоим каждому 
        элементу объединения $\bigcup_{i=1}^\infty A_i$ уникальный 
        натуральный номер. Тем самым построена биекция с множеством $\N$, 
        а значит, объединение счётно.
    \end{quote}
\end{proof}

\subsubsection{Метод диагонали Кантора}

\begin{theorem}[\textbf{Теорема Кантора, 1-я формулировка}]
    Множество $\B^\omega$ всех бесконечных последовательностей из нулей 
    и единиц не является счётным.
\end{theorem}

\begin{proof}[Доказательство методом диагонали Кантора]~\\
    \begin{quote}
        Докажем теорему от противного. Предположим, что множество $\B^\omega$ 
        счётно. Тогда все его элементы можно занумеровать натуральными 
        числами и выписать в бесконечный список строк:
        \begin{align*}
            \alpha_1 &= a_{11} a_{12} a_{13} \dots \\
            \alpha_2 &= a_{21} a_{22} a_{23} \dots \\
            \alpha_3 &= a_{31} a_{32} a_{33} \dots \\
            \vdots \;& \quad \vdots
        \end{align*}
        Рассмотрим бесконечную последовательность элементов, стоящих на 
        главной диагонали таблицы: $\alpha = a_{11} a_{22} a_{33} \dots$.
        Построим новую бинарную последовательность $\beta$, инвертировав 
        каждый диагональный символ ($0 \leftrightarrow 1$):
        \[
            \beta = (1 - a_{11})(1 - a_{22})(1 - a_{33}) \dots
        \]
        По построению $\beta \in \B^\omega$. Так как мы предположили, что в 
        таблице перечислены все элементы $\B^\omega$, должен существовать 
        номер $n \in \N$, для которого $\alpha_n = \beta$.

        Однако для любого $n \in \N$ на $n$-й позиции в строке $\alpha_n$ 
        стоит символ $a_{nn}$, а в последовательности $\beta$ на $n$-й 
        позиции стоит символ $1 - a_{nn} \neq a_{nn}$. Следовательно, 
        $\forall n \in \N : \alpha_n \neq \beta$, и последовательности 
        $\beta$ нет в списке. Полученное противоречие доказывает 
        несчётность множества $\B^\omega$.
    \end{quote}
\end{proof}

\begin{remark}
    Изложенный способ рассуждений называется \textdef{диагональным методом 
    Кантора} и широко применяется в теории множеств, математической логике 
    и теории вычислимости.
\end{remark}

\subsubsection{Булеан и основная теорема Кантора}

\begin{definition}
    \textdef{Булеаном} множества $X$ (обозначается $\mathcal{P}(X)$ или 
    $2^X$) называется множество всех подмножеств множества $X$:
    \[
        \mathcal{P}(X) = \{Y \mid Y \subseteq X\}
    \]
\end{definition}

\begin{example} ~
    \begin{align*}
        \mathcal{P}(\{0, 1, 2\}) &= \{\varnothing, \{0\}, \{1\}, \{2\}, 
        \{0, 1\}, \{0, 2\}, \{1, 2\}, \{0, 1, 2\}\}, \\
        \mathcal{P}(\varnothing) &= \{\varnothing\} \quad 
        \left(\text{при этом } \abs{\varnothing} = 0, \ 
        \abs{\mathcal{P}(\varnothing)} = 1\right).
    \end{align*}
\end{example}

\begin{theorem}[\textbf{Основная теорема Кантора}]
    Никакое множество не равномощно своему булеану:
    \begin{equation} \label{eq:cantor_main}
        \forall X : X \nsim \mathcal{P}(X) \tag{11}
    \end{equation}
\end{theorem}

\begin{corollary}
    Мощность булеана любого множества строго больше мощности самого 
    множества: $\abs{X} < \abs{\mathcal{P}(X)}$.
\end{corollary}

\begin{proof}[Доказательство основной теоремы Кантора]~\\
    \begin{quote}
        Докажем от противного. Пусть существует множество $X$, равномощное 
        своему булеану: $X \sim \mathcal{P}(X)$. Тогда существует биективное 
        отображение $f: X \to \mathcal{P}(X)$.

        Сформируем специальное подмножество $Y \subseteq X$, включив в него 
        те и только те элементы $x \in X$, которые не принадлежат своему 
        образу $f(x)$:
        \[
            Y = \{x \in X \mid x \notin f(x)\}
        \]
        Так как $Y \subseteq X$, то $Y \in \mathcal{P}(X)$. Поскольку $f$ 
        сюръективно (как биекция), у элемента $Y$ обязан существовать 
        прообраз $y \in X$, такой что $f(y) = Y$.

        Выясним, принадлежит ли элемент $y$ множеству $Y$:
        \begin{enumerate}
            \item Пусть $y \in Y$. Тогда по определению множества $Y$ 
            должно выполняться условие $y \notin f(y)$. Но $f(y) = Y$, 
            откуда $y \notin Y$ --- противоречие.
            \item Пусть $y \notin Y$. Так как $f(y) = Y$, это означает, что 
            $y \notin f(y)$. Но тогда по характеристическому свойству 
            множества $Y$ элемент $y$ обязан принадлежать $Y$, то есть 
            $y \in Y$ --- снова противоречие.
        \end{enumerate}
        Оба возможных случая приводят к противоречию (аналогичному 
        парадоксу Рассела). Следовательно, биекции $f$ не существует и 
        $X \nsim \mathcal{P}(X)$.
    \end{quote}
\end{proof}

\subsubsection{Континуум-гипотеза и бесконечные множества}

\begin{remark}[Континуум-гипотеза]
    \textdef{Континуум-гипотеза} Георга Кантора утверждает, что между 
    мощностью счётного множества ($\aleph_0$) и мощностью континуума 
    ($\mathfrak{c} = 2^{\aleph_0}$) не существует промежуточных кардинальных 
    чисел, то есть $\mathfrak{c} = \aleph_1$. Иными словами, любое 
    бесконечное подмножество множества вещественных чисел $\R$ либо счётно, 
    либо равномощно всему $\R$. Работами К.~Гёделя (1940) и П.~Коэна (1963) 
    доказано, что в стандартной аксиоматике теории множеств Цермело---Френкеля 
    с аксиомой выбора (ZFC) континуум-гипотеза неразрешима (независима от 
    аксиом).
\end{remark}

\begin{theorem}
    Любое бесконечное множество содержит счётное подмножество.
\end{theorem}

\begin{proof}[Доказательство о счётном подмножестве]~\\
    \begin{quote}
        Пусть $X$ --- произвольное бесконечное множество. Так как $X$ 
        бесконечно, оно непусто ($X \neq \varnothing$), поэтому в нём можно 
        выбрать некоторый элемент $x_1 \in X$.

        Рассмотрим множество $X \setminus \{x_1\}$. Оно также непусто, 
        иначе множество $X$ состояло бы из единственного элемента $x_1$ и 
        было бы конечным. Выберем элемент $x_2 \in X \setminus \{x_1\}$.

        Продолжая этот процесс по индукции, на $n$-м шаге рассматриваем 
        множество $X \setminus \{x_1, x_2, \dots, x_{n-1}\}$. Оно непусто 
        (в противном случае мощность $X$ равнялась бы $n - 1$), что 
        позволяет выбрать новый элемент 
        $x_n \in X \setminus \{x_1, x_2, \dots, x_{n-1}\}$.

        По аксиоме счётного выбора мы получаем бесконечную 
        последовательность попарно различных элементов 
        $x_1, x_2, \dots, x_n, \dots \in X$. Множество 
        $A = \{x_1, x_2, \dots, x_n, \dots\} \subseteq X$ по построению 
        находится в биекции с $\N$ ($n \leftrightarrow x_n$), а значит, 
        является искомым счётным подмножеством.
    \end{quote}
\end{proof}